%------------------------------------------------------------------------------%
% Luis A. D'Afonseca
%
%------------------------------------------------------------------------------%
\documentclass[a4paper,12pt,fleqn]{article}

\usepackage{style_avaliacao}

\ShowAnswers

%------------------------------------------------------------------------------%
\begin{document}

\makeHeader{Integrais e Séries}{Prova Substitutiva -- Turma 4}{14/12/2023}

\vspace{\baselineskip}
\textbf{Substitui a prova} \underline{\hspace{20mm}}
\vspace{\baselineskip}

Os exercícios 1 e 2 valem 30 para quem perdeu a prova 1 ou 2 \\[-1mm]
O exercício 3 vale 40 para quem perdeu a prova 3 \\[-1mm]
O exercício 4 vale 40 para quem perdeu a prova 4

% 01
%------------------------------------------------------------------------------%
\vspace{\baselineskip}
\exercise{20}
Calcule a integral \(\int_0^3 x^2\ln(x) dx\).
\textit{Atenção ao intervalo de integração}.

\begin{answer}
  Calculando a primitiva
  \(
    F(x) = \int x^2\ln(x) dx
  \),
  por integral por partes
  \[
    u = \ln(x)
    \qquad
    dv = x^2dx
  \]
  \[
    du = \frac{dx}{x}
    \qquad
    v = \frac{x^3}{3}
  \]
  \begin{align*}
    F(x)
    & = \int x^2\ln(x) dx \\[2mm]
    & = \ln(x)\frac{x^3}{3} - \int \frac{x^3}{3} \frac{dx}{x} \\[2mm]
    & = \frac{x^3\ln(x)}{3} - \frac{1}{3}\int x^2 dx\\[2mm]
    & = \frac{x^3\ln(x)}{3} - \frac{1}{3} \frac{x^3}{3} + C\\[2mm]
    & = \frac{x^3}{3^2}\big(3\ln(x)-1\big) + C
  \end{align*}
  Para calcular a integral definida precisamos observar que temos uma assintota
  vertical em $x=0$, portanto, essa é uma integral imprópria
  \begin{align*}
    I
    & = \int_0^3 x^2\ln(x) dx \\[2mm]
    & = \lim_{t\to0} \int_t^3 x^2\ln(x) dx \\[2mm]
    & = \lim_{t\to0} \left[\frac{x^3}{3^2}\big(3\ln(x)-1\big)\evalat{t}{3}\right]  \\[2mm]
    & = \lim_{t\to0} \left[\frac{3^3}{3^2}\big(3\ln(3)-1\big)
                         - \frac{t^3}{3^2}\big(3\ln(t)-1\big)\right]   \\[2mm]
    & = 3\big(3\ln(3)-1\big)
      - \frac{1}{3^2}\lim_{t\to0} \left[t^3\big(3\ln(t)-1\big)\right]
  \end{align*}
  Precisamos calcular o limite
  \begin{align*}
    L
    & = \lim_{t\to0} t^3\big(3\ln(t)-1\big)  \\[2mm]
    & = \lim_{t\to0} \frac{3\ln(t)-1}{t^{-3}} & \text{usando l'Hopital}\\[2mm]
    & = \lim_{t\to0} \frac{3t^{-1}}{-3t^{-4}} \\[2mm]
    & = \lim_{t\to0} -t^3 = 0
  \end{align*}
  Portando
  \[
    I = \int_0^3 x^2\ln(x) dx
    = 3\big(3\ln(3)-1\big)
    = 9\ln(3)-3
  \]
\end{answer}

% 02
%------------------------------------------------------------------------------%
\vspace{\baselineskip}
\exercise{20}
Calcule o volume do sólido construído pela rotação, em torno do eixo $x$, da região
contida entre as curvas \(y=x\) e \(y=x^2\)

\begin{answer}
  Integrando por seções transversais em $x$ temos
  \[
    V = \int_a^b A(x) dx
  \]
  Calculando os pontos de intersecção das curvas
  \begin{align*}
    x & = x^2 \\
    x^2 - x & = 0 \\
    x(x-1) & = 0
  \end{align*}
  temos \(x=0\) ou \(x=1\), então \(a=0\), \(b=1\).
  A área da seção transversal é
  \[
    A(x) = \pi r_1^2 - \pi r_2^2 = \pi \left(x^2-x^4\right)
  \]
  Assim
  \begin{align*}
    V
    & = \int_0^1 \pi \left(x^2-x^4\right) dx \\[2mm]
    & = \pi \int_0^1 x^2-x^4 dx \\[2mm]
    & = \pi \left(\frac{x^3}{3} - \frac{x^5}{5}\right)\evalat{0}{1} \\[2mm]
    & = \pi \left[
          \left(\frac{1^3}{3} - \frac{1^5}{5}\right) -
          \left(\frac{0^3}{3} - \frac{0^5}{5}\right)
        \right] \\[2mm]
    & = \pi \left(\frac{1}{3} - \frac{1}{5}\right)
     = \frac{2\pi}{15}
  \end{align*}
\end{answer}

% 03
%------------------------------------------------------------------------------%
\vspace{\baselineskip}
\exercise{20}
Analise a convergência da série \(\sum_{n=1}^\infty ne^{-n^2}\)

\begin{answer}
  Usando o teste da integral com a \textbf{função real} \(f(x)=xe^{-x^2}\) que é:
  \begin{enumerate}
    \item \textbf{contínua} para todo $x$ real
    \item \textbf{positiva} para $x>0$
    \item para verificar se é \textbf{decrescente} vamos calcular sua derivada
    \begin{align*}
      f'(x)
      & = \frac{d}{dx}xe^{-x^2} \\[2mm]
      & = e^{-x^2} + x\frac{d}{dx}e^{-x^2} \\[2mm]
      & = e^{-x^2} + xe^{-x^2}\frac{d}{dx}\left(-x^2\right) \\[2mm]
      & = e^{-x^2} + xe^{-x^2}(-2x) \\[2mm]
      & = e^{-x^2}\left(1 - 2x^2\right)
    \end{align*}
    A função será decrescente para
    \begin{align*}
      1 - 2 x^2 & < 0                              \\
      x^2       & > \frac{1}{2}                    \\
      x         & > \frac{1}{\sqrt{2}} \approx 0,7
    \end{align*}
  \end{enumerate}
  Temos então que a série converge, se e somente, se a integral
  \[
    \int_1^\infty xe^{-x^2} dx
  \]
  converge.
  Calculando a primitiva por substituição simples
  \[
    u = -x^2
    \qquad
    du = -2xdx
    \qquad
    dx = \frac{du}{-2x}
  \]
  temos
  \begin{align*}
    F(x)
    & = \int xe^{-x^2} dx \\[2mm]
    & = \int xe^u \frac{du}{-2x} \\[2mm]
    & = \frac{-1}{2}\int e^u du \\[2mm]
    & = \frac{-e^u}{2} + C \\[2mm]
    & = \frac{-e^{-x^2}}{2} + C
  \end{align*}
  Calculando a integral imprópria
  \begin{align*}
    I
    & = \int_1^\infty xe^{-x^2} dx \\[2mm]
    & = \lim_{b\to\infty}\int_1^b xe^{-x^2} dx \\[2mm]
    & = \lim_{b\to\infty}\left(\frac{-e^{-x^2}}{2}\right)\evalat{1}{b} \\[2mm]
    & = \frac{-1}{2}\lim_{b\to\infty}\left(e^{-x^2}\right)\evalat{1}{b} \\[2mm]
    & = \frac{-1}{2}\lim_{b\to\infty}\left(e^{-b^2}-e^{-1^2}\right) \\[2mm]
    & = \frac{-1}{2}\left(\lim_{b\to\infty}\left(e^{-b^2}\right)-e^{-1}\right) \\[2mm]
    & = \frac{1}{2e} \in \R
  \end{align*}
  Como a integral converge a série também converge.
\end{answer}

% 04
%------------------------------------------------------------------------------%
\vspace{\baselineskip}
\exercise{20}
Sabendo que \(f^{(n)}(x) = 2^x\ln(2)^n\) para \(n=0, 1, 2, \dots\)
\begin{enumerate}[label=\alph*)]
  \item Construa a Série de Taylor de $f$ centrada em zero
  \item Mostre que a Série de Taylor de $f$ converge para $f$ em $[0, 10]$
\end{enumerate}

\begin{answer}
\begin{enumerate}[label=\alph*)]
  \item A Série de Taylor para essa função é
  \[
    T(x)
    = \sum_{n=0}^\infty \frac{f^{(n)}(0)}{n!}x^n
    = \sum_{n=0}^\infty \frac{2^n\ln(2)^n}{n!}x^n
  \]
  \item Para provarmos que a Série de Taylor converge para a função devemos
  mostrar que \(R_n(x)\to 0\). Pelo Teorema de Taylor temos
  \[
    R_n(x) = \frac{f^{(n+1)}(\xi)}{(n+1)!}x^{n+1}
  \]
  com $\xi$ entre zero e $x$, em nosso caso
  \[
    R_n(x)
    = \frac{2^\xi\ln(2)^n}{(n+1)!}x^{n+1}
    = \frac{\ln(2)^n}{(n+1)!}\, 2^\xi\, x^{n+1}
  \]
  Portanto
  \[
    \abs{R_n(x)}
    = \frac{\ln(2)^n}{(n+1)!}\, 2^\xi\, \abs{x}^{n+1}
  \]
  Como \(0\leq\xi\leq x\leq 10\) temos
  \[
    2^\xi \leq 2^{10}
    \qquad
    \text{e}
    \qquad
    \abs{x}^{n+1}
    = x^{n+1}
    \leq 10^{n+1}
  \]
  Assim
  \[
    \abs{R_n(x)}
    \leq \frac{\ln(2)^n}{(n+1)!}\, 2^{10}\, 10^{n+1}
    = 10\cdot 2^{10} \,\frac{\ln(2)^n\,10^n}{(n+1)!}
    = 10\cdot 2^{10} \,\frac{\big(10\ln(2)\big)^n}{(n+1)!}
  \]
  como
  \[
    \lim_{n\to\infty} 10\cdot 2^{10} \,\frac{\big(10\ln(2)\big)^n}{(n+1)!} = 0
  \]
  o Teorema do Confronto garante que $R_n(x)\to 0$ e portanto a Série de Taylor
  converge para a função.
\end{enumerate}
\end{answer}

%------------------------------------------------------------------------------%
\end{document}
%------------------------------------------------------------------------------%
